\FloatBarrier
\section{FT inner shield}
\label{sec:ft-shield}

The shield/cone interference recorded below is resolved by the subsequent cone revision in
Section~\ref{sec:moller-cone}.

The initial revision below is retained with its original measurements. The final smooth-flange configuration
supersedes it; see Section~\ref{sec:ft-shield-smooth}.

The ninth beamline study renames taggerInnerShield to FTInnerShield in all nine source and runtime CAD
variations that contain it and in nine GEMC2 SQLite CAD rows. The tungsten-alloy material beamline\_W,
color, mother, and original passive placement are retained: z = 1643.00 mm with a 180-degree y rotation.
This downstream grooved tube is distinct from the previously optimized FTInnerShieldAndFlange at z = 1273.27
mm.

The original contains 44,122 triangles, 7,404 boundary edges, 67 surface components, 351 non-manifold
vertices, and 1,750 faces selected by the intersection check. Its open and non-manifold surface has no
defined solid volume. Routine welding and generic hole filling do not establish a sound surface. A
deterministic reconstruction uses the measured tube cross-section, chamfers, flange fillets, and feature
positions, without Poisson reconstruction. The central bore remains 60 mm in diameter, nominal tube outer
diameter 80 mm, length 665 mm, and flange outer radius 100 mm. The flange occupies local z = -625 to -605
mm.

Two longitudinal grooves at local azimuths 45 and 135 degrees are retained after the flange. Each is 6 mm
wide and 3 mm deep, extending from local z = -575 mm to the tube end at z = 0. A 6 mm axial end radius
blends its floor into the tube. All 26 flange through-holes and the single radial tube hole are closed. The
requested 0.5 mm recess threshold is applied to axial depth: eight 5.75 mm deep conical recesses, sixteen
0.6 mm face chamfers, and two 10 mm deep counterbores remain as blind features. The conical mouths
originally touch the flange outer rim exactly; reducing each mouth radius from 10.00 to 9.95 mm removes that
tangency while retaining their depths and 4.25 mm bottom radii.

The reference uses 128 angular segments for the main profile, sampled fillet/end curves, and Boolean cuts
for the grooves and retained recesses. A 0.000001 mm weld removes near-coincident Boolean vertices and their
sliver facets before saving the reference. Its 15,100 triangles define a closed manifold component of genus
one with volume 1938424.030371 mm$^3$. Protected quadric simplification retains 9,404 groove, recess, and
corner faces at their saved float32 coordinates and produces 13,500 triangles, a 69.40\% reduction from the
engineering mesh. File size decreases from 2,206,184 to 675,084 bytes.

The final mesh remains one closed manifold component of genus one, with no selected self-intersections or
zero-area facets. Bounds remain [-100, -100, -665] to [100, 100, 0] mm apart from negligible source z
rounding. Minimum triangle altitude is 0.000344255 mm. Bidirectional sampling against the rebuilt reference
gives a maximum change of 0.023925 mm, below the 0.05 mm guard. Simplification changes reconstructed volume
by -0.00058479\%. These are sampled checks, not proven global surface bounds or a comparison to a certified
original mass.

Independent VTK containment probes verify all 58 requested inside/outside results across groove floors,
blind recesses, filled bolt shafts, the filled tube hole, and the central passage. Outside the removed
drilling walls and flange drilling region, retained source profile vertices lie within 0.050719 mm of the
reconstructed surface. One-way vertex, edge, and face sampling of the retained raw tube/rim subset gives a
larger maximum 0.969091 mm: the reconstruction replaces coarse, disconnected raw facets with the measured
nominal profile. This original-to-final difference is distinct from the much smaller subsequent decimation
error. The original mass change and detector-response equivalence are not established. The builder
reproduces the final STL and reconstructed reference byte for byte from the archived engineering input.

\begin{table}[htbp]
\centering
\caption{Original and optimized downstream FT inner shield.}
\label{tab:ft-shield-mesh}
\begin{tabular}{lrr}
\toprule
Quantity & Original & Optimized \\
\midrule
Triangles & 44,122 & 13,500 \\
Bytes & 2,206,184 & 675,084 \\
Boundary edges & 7,404 & 0 \\
Surface components & 67 & 1 \\
Solid volume (mm$^3$) & Undefined & 1938412.695 \\
\bottomrule
\end{tabular}
\end{table}

A two-volume assembly includes the unchanged main-cone at z = 1589.27 mm and the shield at its original
placement. The 10,000-point Geant4 test reports five shield/cone overlaps in each case, confirming the
pre-existing interference remains. Original construction has three solid warnings; optimized assembly
construction has one warning from the unchanged cone. The isolated optimized shield loads with no Geant4
exceptions. Full default beamline construction completes without missing meshes, with 1 solid warning and 9
overlap reports involving the shield/cone and downstream/adapter geometry. The original placements are
retained rather than adjusted in this study. All fifteen native geometry/material and fourteen CAD
variations pass the semantic comparison against GEMC2.

Two independent tracking comparisons isolate the shield in air. The first uses the defective engineering
mesh as baseline; the second uses the hole-filled, grooved, unreduced reference to measure simplification
after repair. Each uses the fixed GEMC3 executable snapshot from the flange study, matched 5 GeV
target-origin primaries uniform in solid angle over 0.7--2.7 degrees and full azimuth, FTFP\_BERT without
optics or magnetic field, one thread, default production cuts, and no GUI or event output. Five
alternating-order repetitions use seeds 12345--12349 and 200 events per electron, proton, positive-pion, and
neutron run. Each comparison contains 52 processes, including zero-event startup and isolated overlap
checks. All 104 processes complete normally and all eighty transport runs report 200 events.

Engineering-original/optimized median transport wall-time ratios are 0.923, 1.007, 0.992, 0.924 for
electrons, protons, positive pions, and neutrons respectively. Ratios below one denote slower optimized
tracking. The engineering surface carries solid and navigation warnings; repair and authorized hole closure
change material boundaries, so this comparison does not attribute timing changes solely to triangle count.
Original material volume is undefined.

\begin{table}[htbp]
\centering
\caption{Median transport wall seconds per 200 events: engineering comparison.}
\label{tab:ft-shield-engineering}
\begin{tabular}{lrrr}
\toprule
Particle & Original & Optimized & Original/optimized \\
\midrule
e- & 22.46020 & 24.33720 & 0.923 \\
proton & 2.23604 & 2.22142 & 1.007 \\
pi+ & 2.96617 & 2.99009 & 0.992 \\
neutron & 1.99370 & 2.15824 & 0.924 \\
\bottomrule
\end{tabular}
\end{table}

Rebuilt-reference/optimized median ratios are 1.087, 1.075, 1.075, 1.049 in the same particle order. This
comparison retains the same blind recesses, closed holes, grooves, and material, with the recorded small
decimation volume difference. Both isolated reference and optimized fixtures have no Geant4 exceptions.
These short workload results are reported separately from the repair comparison; small ratios are sensitive
to timing variation and stochastic transport. They do not establish full CLAS12 speedup or response
equivalence. The optimized mesh is faster in all 20 matched particle/seed pairs in the reference comparison.
Median startup simulation time changes from 1.063410 to 0.770623 s, with peak startup RSS 199.94 to 136.06
MiB.

\begin{table}[htbp]
\centering
\caption{Median transport wall seconds per 200 events: reconstructed-reference comparison.}
\label{tab:ft-shield-reference}
\begin{tabular}{lrrr}
\toprule
Particle & Reference & Optimized & Reference/optimized \\
\midrule
e- & 26.31760 & 24.21180 & 1.087 \\
proton & 2.37224 & 2.20600 & 1.075 \\
pi+ & 3.16977 & 2.94966 & 1.075 \\
neutron & 2.28295 & 2.17616 & 1.049 \\
\bottomrule
\end{tabular}
\end{table}

\begin{table}[htbp]
\centering
\caption{Geant4 exception counts across each isolated 52-process comparison.}
\label{tab:ft-shield-warnings}
\begin{tabular}{lrrr}
\toprule
Diagnostic & Original & Optimized & Reference \\
\midrule
GeomSolids1001 & 52 & 0 & 0 \\
GeomSolids1002 & 0 & 0 & 0 \\
GeomNav1002 & 45 & 0 & 0 \\
GeomNav0003 & 0 & 0 & 0 \\
GeomVol1002 & 0 & 0 & 0 \\
\bottomrule
\end{tabular}
\end{table}

\FloatBarrier
\begin{figure}[htbp]
\centering
\includegraphics[width=0.9\textwidth]{figures/ft_shield_stl_comparison.pdf}
\caption{Original FT inner shield (left) and rebuilt, simplified shield (right). Matched cameras and bounds
use original coordinates; longitudinal grooves are retained.}
\label{fig:ft-shield-stl-comparison}
\end{figure}

\FloatBarrier
\begin{figure}[htbp]
\centering
\includegraphics[width=0.9\textwidth]{figures/ft_shield_details.pdf}
\caption{Matched flange underside views and tube sections at local z = -100 mm. Flange recesses remain
blind, and both longitudinal grooves remain after through-hole closure.}
\label{fig:ft-shield-details}
\end{figure}
